% titel: Blatt 10, A3: Brønsted- und Lewis-Säuren/Basen
% skript: Mortimer Kap. 17 (Säure-Base-Begriffe)

% Teil I: Werkzeuge, kein Lösungsweg.

Zwei Deutungsebenen: Brønsted (Protonen-Transfer) und Lewis
(Elektronenpaar-Akzeptor/Donator).

\begin{werkzeugkasten}[Brønsted vs.\ Lewis]
\begin{itemize}
  \item \textbf{Brønsted-Säure:} gibt H$^+$ ab (H enthaltend, z.~B. HNO$_3$,
        HCN).
  \item \textbf{Brønsted-Base:} nimmt H$^+$ auf (Anion oder Molekül mit
        freiem Elektronenpaar an elektronegativem Atom).
  \item \textbf{Lewis-Säure:} Elektronenpaar-\emph{Akzeptor} (oft
        ungesättigte Zentralatome, z.~B. BF$_3$, AlCl$_3$).
  \item \textbf{Lewis-Base:} Elektronenpaar-\emph{Donator} (freies Paar an
        O, N, S...).
\end{itemize}
Eine Linie (HCN, HNO$_3$) kann beides sein (H$^+$- und Elektronenpaar-Rolle
je nach Partner). Es gibt auch reine Lewis-Säuren ohne H (AlCl$_3$, BF$_3$).
\textit{Skript:} Mortimer Kap. 17 (Säure-Base-Theorien).
\end{werkzeugkasten}

\begin{achtung}[Darauf achten]
\begin{itemize}
  \item Enthält ein Molekül \emph{kein} abstrahierbares H (BF$_3$,
        AlCl$_3$), kommt nur Lewis-Säure in Frage.
  \item mit freiem Elektronenpaar (Diethylether, Aceton, tert-Butanolat-Ion)
        sind Lewis-Basen; können dazu Brønsted-Base sein.
  \item Ionen wie Na$_2$CO$_3$ liefert das Carbonat CO$_3^{2-}$ als Base.
\end{itemize}
\end{achtung}
